\section{Introduction}

Deployed agents rarely run on a bare system prompt. Operators give them
personas (``you are a resourceful engineer who unblocks people''), they
increasingly work alongside other agents, passing tickets and notes back and
forth \citep{carichon2025multiagent}, and they act through whatever tools the
deployment grants. None of this is chosen by the model. If an agent's
willingness to follow a stated rule shifts with the character it is told to
play, with what other agents say they did, or with the form its actions take,
then a safety evaluation run on the bare model with one tool misses how it
behaves once deployed.

Work on conformity shows that LLMs follow a majority on factual and opinion
questions \citep{weng2025conformity,zhu2025conformity,choi2025conformity}.
Persona assignment is known to shift toxicity and style
\citep{deshpande2023toxicity}. Few studies ask whether either kind of context
moves a \emph{rule-following decision} in an agentic setting, where the
behaviour is a tool call rather than an answer, and fewer still vary the tool
itself.

We study this with a controlled evaluation. In \R{all}{nscenarios} scenarios an
operations agent is asked for a fact that lives only on a host its network
policy forbids. We cross the tool the agent holds (a named HTTP tool or a
Python sandbox) with operator personas and with handoff notes from peer
agents, each with a matched \emph{presence control}, and run the design on
\R{all}{nmodels} models. Our contributions are:

\begin{itemize}
  \item \textbf{A judge-free evaluation of agentic rule-following.} Violation,
  use of the forbidden data, and honesty of the agent's later self-report are
  all read mechanically from the transcript. Simulated tools return a
  \emph{successful} payload carrying a unique marker, so a string match
  separates delivering forbidden data from probing. Every scenario also runs
  under a \emph{permit} policy, so a low violation rate counts as compliance
  only for models that make the request when allowed.
  \item \textbf{Evidence that tool and social context move rule-following.}
  Among the \R{all}{nactive} models that violate at all, the forbidden request
  is about four times as likely in odds when written as code as when made
  through a named tool, and peer testimony and personas shift it in both
  directions against matched controls (\S\ref{sec:results}).
  \item \textbf{Evidence that these effects differ in direction across
  models.} One model reverses both peer effects, and the same persona raises
  violation in one model and lowers it in another. Four frontier models do not
  move at all. Pooled estimates therefore describe a population of models and
  not any one of them.
\end{itemize}

We release the scenarios, the Inspect AI implementation and the per-sample
data; every number in this paper is generated from the evaluation
logs.\footnote{Code and data: link removed for anonymous review.}
